\section{Introduction} \label{sec:intro}
 We focus on quantifying the difficulty of fixed-budget best-arm identification (BAI) problem in non-stationary linear bandits. The linear bandit framework generalizes the classical multi-armed bandit problem by associating each arm with a feature vector and modeling rewards as linear in an unknown parameter. Algorithms such as UCB and Thompson Sampling are known to perform optimally in the well-studied regret-minimization setting \citep{lattimore2020bandit}. However, in the best-arm identification setting, where exploration is more important than cumulative reward, these algorithms are suboptimal \citep{bubeck2009pure}, thus requiring specialized attention. In the standard BAI setting, a stationary environment is assumed where the rewards of arms are sampled i.i.d \citep{audibert2010best, soare2014best, NEURIPS2019_8ba6c657}. However, algorithms designed for these environments can completely fail as soon as this assumption is lifted, for example, in settings where the value of an arm can change at any time step \citep{abbasi2018best,xiong2024b}. An interesting question is \emph{exactly how the difficulty of this setting changes once this stationary assumption is lifted}. 

In recent work, \citet{xiong2024b} partially answered this question by showing that when the arm set is restricted to the standard basis, the difficulty of non-stationary BAI scales proportionally with the dimension. While informative, this result is ultimately unsatisfying as it collapses the linear bandit model back into a multi-armed bandit. Basis arm sets erase correlations between arms, precisely the geometric structure that linear bandits are meant to exploit. With this in mind, we aim to uncover precisely which relationship between arms truly governs the difficulty of non-stationary BAI. Inspired by \citet{soare2014best}, a natural conjecture is that the pairwise relationships between every arm should be considered when measuring difficulty. However, such a characterization implies that any arm set containing a basis is no easier than the basis itself, and thus no progress is made. To refine this view, we introduce adjacency, a geometric notion that captures which pairs of arms can compete for optimality. Accordingly, we find that the pairwise relationships between adjacent arms are necessary and sufficient to measure the difficulty of non-stationary BAI. Formally, by leveraging the adjacent structure of the arm set, we refine the previous notion of difficulty in non-stationary BAI by establishing a measure of complexity that adapts to the geometry of \textit{any} arm set, which we term as \textit{arm-set-dependent}.


\paragraph{Our contributions.}
\begin{itemize}
    \item \textbf{Adjacency and Non-Stationary BAI.} The main novelty of our paper stems from our characterization of non-stationary BAI as a problem of differentiating only between adjacent arms, a notion we formalize in later sections. The intuition for this restriction stems from our central Lemma \ref{lem:adjacency_optimal}, which implies that if an arm is better than its adjacent arms, then it is the optimal arm. Accordingly, we introduce a new arm-set-dependent complexity measure $ H_{\mathrm{Adjacent}}(\X)$ and show it strictly refines the minimax-optimal complexity $ H_G$. 
    \item \textbf{Arm-set-dependent lower bound.} We show in Theorem \ref{thm:finallower} that for fixed arm set $\X$, the error probability of any algorithm must be at least $\exp\left(-O\left(T/H_{\mathrm{Adjacent}}(\X)\right)\right)$. Our lower bound builds off of previous works \citep{kaufmann2016complexity, abbasi2018best} concerning lower bounds in best-arm identification for the multi-armed bandit setting. However, these methods alone are not sufficient for the linear setting, which we address by characterizing the lower bound as an optimization problem in Lemma \ref{thm:optlower}, and which we solve in Lemma \ref{lem:equality} by exploiting the adjacent geometry of the arm set.
    \item \textbf{Matching upper bound.} We introduce the Adjacent-optimal design, a specialization of the well-known $\X\Y$-optimal design \citep{soare2014best}, that instead reduces variance only between adjacent arms. Using the Adjacent-optimal design, we introduce the algorithm \algoname. Utilizing Lemma \ref{lem:adjacency_optimal}, we obtain Theorem \ref{thm:upper_bound1} showing that for fixed arm set $\X$, the error probability of \algoname is at most $\exp\left(-\Omega\left(T/H_{\mathrm{Adjacent}}(\X)\right)\right)$, verifying the tightness of our lower bound, and establishing $ H_{\mathrm{Adjacent}}(\X)$ as 
    the arm-set-dependent complexity measure of this setting.
\end{itemize}
%\textcolor{magenta}{Contributions read much better now, thanks!}

